\section{Discussion and limitations}
\label{sec:discussion}

\paragraph{What the evidence says.}
``How much investment is needed'' has no single answer: it is a distribution over the target's rank, the aggregation rule, the cap and the cost measure.
Three statements hold, each with its conditions; the ordering of rules by cost is the least robust of them.
(1) Under the permissive cap (ceiling) and normalised scores, the mean-like rules (arithmetic, geometric, harmonic) cost the most points. The cheapest rules are median, Borda and min on OLLB (3.6, 6.3, 10.4 points for the 20 closest challengers; plurality 16.4 is above them) and plurality, min and median on HELM; Borda and the Condorcet requirement lie in between and are strongly dataset-dependent (Section~\ref{sec:exp-rules}). This ordering does \emph{not} survive every check: on raw OLLB scores the mean is no longer the most expensive rule ($\tau=0.21$ against the main ordering), and without MuSR the ordering is reshuffled almost completely ($\tau=-0.07$); it holds on the no-merge and top-500 slices (Section~\ref{sec:exp-robust}).
(2) Under a cap that limits the raise (fixed $+10$ points, or the empirical cap from fine-tuning gains), whether first place is reachable at all depends on rank and rule, and on OLLB the mean is the \emph{most} reachable rule while min, median and Condorcet are the least; on HELM the order differs (Section~\ref{sec:exp-caps}). A rule that is cheapest to win under one cap is therefore not the hardest to win under another, and no rule is uniformly ``harder to game'' in our data.
(3) Targeting a few tasks saves points over a uniform improvement under every rule except the mean, where it only concentrates the same points (Section~\ref{sec:exp-uniform}).
The contradicted expectation is instructive: order-statistic rules do not demand more points than the mean, only specific tasks, and their cost for a single target is the least stable under noise.

\paragraph{Where the conclusions fail.}
The ordering of rules fails under raw scores and without MuSR (above). With four targets in ranks 2--5 we say nothing about the closest challengers beyond ``often indistinguishable from the mean''. Where the cap binds, medians are over survivors and slopes are biased flat.

\paragraph{Limitations.}
\emph{Model.} One target acts alone; rivals do not react, so results are minimum costs under a static leaderboard, not predictions of developer behaviour or equilibria. Raising a score on one task is assumed to cost nothing on others and the score is assumed controllable to any value up to the cap.
\emph{Caps.} The caps are the central assumption. Our empirical cap is the 95th/99th percentile of observed fine-tune-minus-base gains on 540 OLLB pairs; this is an indirect proxy for what contamination or targeted data could achieve, not a measurement of it. HELM has no such pairs, so its empirical cap is a transfer of one OLLB task's quantile.
\emph{Data.} Two leaderboards of one benchmark family, one of them with only six tasks (so median and trimmed rules are coarse), and many OLLB entries are merges or fine-tunes of one another; intervals are over targets, not over leaderboards.
\emph{Coverage.} On OLLB, Borda, Condorcet, plurality and the trimmed mean were solved on a 57-target sample, the trimmed mean has no $K$, Copeland was solved only on HELM, the convex cost only for mean and geometric mean, and weighted means only for the mean. The robustness draws use ten (noise, task bootstrap) or two (task subsets) challengers and a single cap.
\emph{Compute.} The experiments exceed the task's 4 CPU-hour estimate; MILP time-limit flags are not recorded for the robustness draws (Appendix~\ref{app:repro}).

\paragraph{Outlook.}
Extensions: reacting rivals (cf.\ \citealt{chen2026leaderboard}), caps measured from contamination experiments, leaderboards with more tasks.
